\documentclass{subfiles}
\begin{document}
    As stated previously, Huffman decoding needs the tree used in compression 
        to correctly decode the bit-stream; it's obvious that, the larger the 
        source alphabet is, the larger is the outputed tree.
        Can we avoid transmitting the whole tree then?
        In the present section we describe a variation of the classical Huffman 
        encoding that does exactly so. We'll show that all that we need are 
        the symbols to transmit, ordered non-increasingly, and the first 
        codeword for each length.

    Before we proceed, let us observe something: there exist optimal codes 
        that correspond to no Huffman tree, for instance the tree 
        on the right in \Cref{fig:nonHuffman}, corresponds to no Huffman encoding,
        yet it's optimal.
        
    Let's now consider how to compute the Canonical Huffman encoding.
        Starting from the Huffman encoding of the source,
        compute the length of the codewords and construct the following data structures:
        \begin{itemize}
            \item A \lstinline{num} array, to store at \lstinline{num[l]} the number 
                of codewords of length \(l\).
            \item A \lstinline{symb} array, that at \lstinline{symb[l]} stores 
                the list of symbols having codeword length of \(l\)-bits.
            \item A \lstinline{fc} array, storing at \lstinline{fc[l]}
                the first codeword off all symbols encoded with \(l\)-bits.
        \end{itemize}

        \subfile{../../TikZ/Figure 3 - Example of optimal tree not produced by Huffman}
        
        Using the above structures, we assign consecutive codewords to the symbols 
            in \lstinline{symb[l]} starting from \lstinline{fc[l]}. 
            All the arrays are 1-based.

    Building the \lstinline{num} and \lstinline{symb} array is trivial,
        the question is how to build \lstinline{fc}. Let \lstinline{MAX}
        be the maximal codeword length, then \lstinline{fc[MAX] = 0},
        as a consequence of the tree we are building, for any \(l \le\) \lstinline{MAX},
        \lstinline{fc[l] = fc[l + 1] + num[l + 1] / d}, where \(d\) is the size of the alphabet.
        The pseudo-code in \autoref{proc:pseudoFC} uses a bottom-up approach, hence the tree built is 
        unbalanced to the left.
        
        \subfile{../../TikZ/Procedure - FC procedure}

        \begin{example*}
            Assume we want to compute the Canonical Huffman encoding for the tree 
                in \Cref{fig:exampleHuffman}. It follows easely that \lstinline{num = \{1, 6\}}
                and \lstinline{symb = \{a, \{g, b, c, e, f, d\}\}}. We now proceed building the 
                \lstinline{fc} array. By using the procedure we have \lstinline{fc[2] = 0}
                    while \lstinline{fc[1] = (fc[2] + num[2]) / 3 = 2}. The tree we obtain, 
                    symmetric to that of \Cref{fig:exampleHuffman}, is shown below.
                    \subfile{../../TikZ/Figure - Example of Canonical Huffman encoding}
        \end{example*}

    To conclude our discussion, we provide the Procedure \ref{proc:decodeCanonical}
        to decode string encoded using Canonical Huffman.
        Recall that the decoder has access to both \lstinline{fc}
        and \lstinline{symb}; therefore, we can simply read bits until the value 
        is less than \lstinline{fc[l]}, for all \(l\).
        \subfile{../../TikZ/Procedure - Canonical Huffman decoding}
\end{document}
